\subsection{\texorpdfstring{The spaces \(E_{A}\) (A bounded)}{The spaces E\_\{A\} (A bounded)}}

Let E be a locally convex space and A be a convex balanced subset of E. We recall that \(E_{A}\) denotes the vector space generated by A, with \(p_{A}\) the gauge of A, as the semi-norm (II, p.~26, Example 3). We verify immediately that the canonical injection of \(E_{A}\) into E is continuous if and only if A is bounded. If, in addition, E is Hausdorff, a bounded set A does not contain a line (III, p.~2, Remark 4) and so \(p_{A}\) is a norm (II, loc. cit.).

We shall say that a uniform space X is semi-complete if every Cauchy sequence in X is convergent. A complete uniform space is semi-complete; but the converse is not always true (GT, II, § 4, exerc. 4); however, a metrizable semi-complete space is complete (GT, IX, § 2, No.~6, prop. 9).

PROPOSITION 8. --- Let E be a Hausdorff locally convex space and let A be a closed, balanced, bounded and convex subset of E. Let \((x_{n})\) be a Cauchy sequence in \(E_{A}\) . Then this sequence converges in \(E_{A}\) if and only if it converges in E.

The canonical injection from \(E_{A}\) into E is continuous. Hence, if \((x_{n})\) converges in \(E_{A}\) , it converges in E. Conversely, suppose \((x_{n})\) converges to y in E. There exists an increasing sequence of integers \((n_{k})\) such that \(p_{\mathbf{A}}(x_{m}-x_{n}) \leqslant 2^{-k-1}\) if \(m \geqslant n_{k}\) and \(n \geqslant n_{k}\) . Therefore the sequence \((x_{n_{k}} + 2^{-k}\mathbf{A})\) is decreasing. Since A is closed in E, we have \(y \in \bigcap_{k} (x_{n_{k}} + 2^{-k}\mathbf{A})\) , which proves that \((x_{n_{k}})\) , hence \((x_{n})\) , converges to y in \(E_{A}\) .

COROLLARY. --- If A is semi-complete (in particular, complete) then \(E_{A}\) is a Banach space.

In fact, a Cauchy sequence in \(E_{A}\) is also a Cauchy sequence for the topology of E and is contained in a set homothetic to A, hence converges in E.

